% ------------------------------------------------------------
% §7  Computability, complexity, and barriers
% ------------------------------------------------------------
\section{The view from computability}\label{sec:computability}

The gap of \cref{sec:gap} is a statement about missing \emph{positive}
knowledge. It is worth closing the overview by asking the complementary,
sharper question: what is \emph{provably} true about the difficulty of
breaking \SHA{}? Here the news is a genuine mix of hard theorems and hard
barriers, and the boundary between them is itself part of the object's
significance.

\subsection{Breaking \SHA{} is an \textsc{np} search problem}
\label{sec:npsearch}

Fix a target digest $y$ and ask for a preimage: some $x$ with
$\SHA(x) = y$. Because \SHA{} is computable by a small circuit
(\cref{sec:history}), a candidate $x$ can be \emph{checked} in linear
time, so preimage-finding is an \textsc{np} search problem in the exact
sense of Cook and Levin~\citep{CookLevin1971}: the ``inversion''
relation is polynomial-time decidable, and the only question is the cost
of \emph{search}. Encode one evaluation of the compression function as a
Boolean circuit---about $2^{16}$ gates---and the equation
$\SHA(x)=y$ becomes a concrete instance of the satisfiability problem
\textsc{sat}, the canonical \textsc{np}-complete language.

The translation is worth seeing once, because it is almost embarrassingly
direct (\cref{fig:tseitin}). Give every wire of the circuit its own
Boolean variable; each gate then constrains its three wires by a handful
of clauses, satisfied exactly by the gate's truth table; pin the $256$
output variables to the bits of $y$. The conjunction of all clauses is
satisfiable \emph{if and only if} a preimage exists, and any satisfying
assignment, read off the input wires, \emph{is} a preimage. The whole of
\SHA{}'s round structure---$\Sigma$'s, carries, message schedule---flattens
into a few hundred thousand such clauses, syntactically indistinguishable
from any other \textsc{sat} instance. A modern conflict-driven
(``\textsc{cdcl}'') solver then does something almost mathematical with
it: it explores partial assignments, and each dead end is analyzed into
a \emph{learned clause}---a small lemma, added to the formula, recording
why that region of the search space is empty. A solver run is a
proof-by-cases of astronomical width, conducted by a machine that
invents its own case analysis.

\begin{figure}[ht]
  \centering
  \begin{tikzpicture}[font=\small,
      wire/.style={thick},
      lbl/.style={font=\scriptsize}]
    % --- little AND gate ---
    \draw[wire] (-1.0,0.55) node[left, lbl] {$u$} -- (0,0.55);
    \draw[wire] (-1.0,-0.15) node[left, lbl] {$v$} -- (0,-0.15);
    % gate body (D shape)
    \draw[thick, fill=projfill]
      (0,-0.5) -- (0.55,-0.5) arc[start angle=-90, end angle=90,
      radius=0.7] -- (0,0.9) -- cycle;
    \node at (0.42,0.2) {$\wedge$};
    \draw[wire] (1.25,0.2) -- (2.1,0.2) node[right, lbl] {$w$};
    % --- arrow to clauses ---
    \draw[->, thick] (2.9,0.2) -- (3.6,0.2);
    \node[anchor=west, align=left, font=\small] at (3.75,0.2)
      {$(\lnot u \vee \lnot v \vee w)$\\
       $(u \vee \lnot w)$\\
       $(v \vee \lnot w)$};
    % --- caption line under both ---
    \node[lbl, align=center, anchor=north] at (3.1,-1.0)
      {one fresh variable per wire, a few clauses per gate;\\
       ${\sim}2^{16}$ gates $\Rightarrow$ a few $\times 10^{5}$ clauses;
       pin $256$ output literals to $y$ and solve.};
  \end{tikzpicture}
  \caption{The reduction that makes ``invert \SHA{}'' a \textsc{sat}
  instance: each gate becomes clauses that are satisfied exactly by its
  truth table (shown: $w = u \wedge v$). Nothing about the construction
  is clever, and that is the point---membership in \textsc{np} is a
  statement about \emph{checkability}, and it holds for \SHA{} as
  mechanically as for any puzzle whose solutions can be verified.}
  \label{fig:tseitin}
\end{figure}

This is not a metaphor: the reduction is run in practice, and by now the
solvers are genuine instruments of record. As early as 2006, Mironov and
Zhang reproduced Wang-style MD5 collisions with an off-the-shelf
\textsc{sat} solver, no differential expertise wired in---the solver
rediscovering, clause by learned clause, constraints it had never been
told~\citep{MironovZhang2006}. Today the traffic runs both directions:
the current record semi-free-start collision on $39$-round \SHA{} was
found with a \textsc{sat}/\textsc{smt} search for the differential
trails of \cref{sec:craft}~\citep{LiLiuWang2024}, and a programmatic
\textsc{sat} solver---one that injects cryptanalytic knowledge as it
runs---produced the first practical $38$-round semi-free-start collision
independently of the classical toolchain~\citep{AlamgirNejatiBright2024}.
The empirical shape of all this effort is a cliff: instances for few
rounds collapse in milliseconds, then somewhere past round $20$ the
running times turn vertical, and no amount of solver engineering has
moved the wall more than a round or two per decade. The cliff's very
existence is one more coordinate in the natural history of the
function---and, unlike most such coordinates, it is cheap to measure.

\begin{project}{The \textsc{sat} cliff}{programming; propositional
  logic}{4--6 weeks}
  Build (or generate with existing tools) \textsc{cnf} encodings of
  $r$-round \SHA{} preimage instances (\cref{fig:tseitin}) and feed
  them to two or three modern solvers. Chart solving time against $r$
  until timeout; then vary what is pinned (full digest, partial digest,
  collision instead of preimage) and watch the cliff move. Plot, on the
  same axis, where the statistical batteries start failing
  (\cref{sec:gap}'s project) and where the differential trails of
  \cref{sec:siege} die: three instruments' frontiers, never to our
  knowledge drawn on one chart. The deliverable is that chart---a
  measured portrait of ``hardness turning on'' round by round.
\end{project}

But note precisely what \textsc{np}-completeness does and does not say.
It says the \emph{worst case} of \textsc{sat} is hard unless
$\mathrm{P}=\mathrm{NP}$; it says almost nothing about the specific,
highly structured instances that hashing produces. Cryptography needs
\emph{average-case} and \emph{one-way} hardness---hard instances that are
easy to \emph{generate} with a solution---and this is a strictly stronger
and more delicate demand than worst-case \textsc{np}-hardness. No
reduction shows that inverting \SHA{} is \textsc{np}-hard, and this is
not merely an unfilled cell in the table: there is a theorem-shaped
reason to expect no such reduction. Akavia, Goldreich, Goldwasser, and
Moshkovitz showed that a black-box reduction (of the standard,
non-adaptive kind) from an \textsc{np}-hard problem to \emph{inverting a
one-way function} would itself trigger complexity-theoretic collapses
nobody believes~\citep{AkaviaGGM2006}. The missing bridge between
$\mathrm{P} \ne \mathrm{NP}$ and \SHA{}'s security is not awaiting a
clever graduate student; by the best current evidence, the standard
bridge-building technology provably cannot span that river.

\subsection{Five worlds, one function}\label{sec:worlds}

The distinction just drawn---worst-case hardness versus generatable
hardness---organizes the deepest open question in the area, and
Impagliazzo gave it an unforgettable form: five possible worlds, exactly
one of which is real~\citep{Impagliazzo1995}. \Cref{fig:worlds} draws
the ladder and scores our function in each world. The point of the
exercise for this paper: \emph{whether an object like \SHA{} can exist
at all} is not a detail of cryptography but the exact content of the
rung between Pessiland and Minicrypt---the existence of one-way
functions---and it is wide open. Everything this paper treats as
mysterious about one specific function is, one level up, the mystery on
which the classification of computational reality turns.

\begin{figure}[ht]
  \centering
  \begin{tikzpicture}[font=\small,
      wname/.style={font=\bfseries\small, anchor=west},
      wdesc/.style={font=\scriptsize, align=left, anchor=west,
                    text width=82mm, inner sep=0pt},
      lbl/.style={font=\scriptsize}]
    % strata, bottom = algorithmica; each row 1.2cm tall, 11.8cm wide
    \fill[padfill]  (0,4.8) rectangle (11.8,6.0);
    \fill[projfill] (0,3.6) rectangle (11.8,4.8);
    \fill[lenfill]  (0,2.4) rectangle (11.8,3.6);
    \fill[red!7]    (0,1.2) rectangle (11.8,2.4);
    \fill[red!14]   (0,0)   rectangle (11.8,1.2);
    \draw (0,0) rectangle (11.8,6.0);
    \foreach \y in {1.2,2.4,3.6,4.8} \draw (0,\y) -- (11.8,\y);
    \node[wname] at (0.2,5.4) {Cryptomania};
    \node[wdesc] at (2.85,5.4)
      {trapdoors exist: public-key cryptography, the internet as
       deployed (\SHA{} still does the hashing)};
    \node[wname] at (0.2,4.2) {Minicrypt};
    \node[wdesc] at (2.85,4.2)
      {one-way functions exist: \SHA{} can be what it appears---PRGs,
       signatures, proof-of-work all sound};
    \node[wname] at (0.2,3.0) {Pessiland};
    \node[wdesc] at (2.85,3.0)
      {hard random instances abound, but none can be generated
       \emph{with its solution}: no one-wayness; \SHA{} invertible};
    \node[wname] at (0.2,1.8) {Heuristica};
    \node[wdesc] at (2.85,1.8)
      {\textsc{np} hard only in the worst case; typical instances
       easy---\SHA{} broken in practice};
    \node[wname] at (0.2,0.6) {Algorithmica};
    \node[wdesc] at (2.85,0.6)
      {$\mathrm{P}=\mathrm{NP}$ (in spirit): every preimage findable
       fast; cryptography impossible, algorithmics paradise};
    % the one-way-function line
    \draw[very thick, dashed, citewine] (0,3.6) -- (12.35,3.6);
    \node[lbl, citewine, anchor=north west, align=left, text width=21mm]
      at (12.0,3.45) {the one-way-\\function line};
    % brace: SHA's conjectured home
    \draw[decorate, decoration={brace, amplitude=6pt}]
      (12.0,3.7) -- (12.0,5.95)
      node[midway, right=8pt, lbl, align=left, text width=21mm]
      {if \SHA{} is what it appears to be, the real world is here};
  \end{tikzpicture}
  \caption{Impagliazzo's five worlds: mutually
  exclusive possibilities for how computation works, exactly one of
  them actual, none of them ruled out. \SHA{}'s conjectured properties
  assert that we live above the one-way-function line (at least
  Minicrypt)---a claim strictly \emph{weaker} than what public-key
  cryptography bets on, and strictly \emph{stronger} than
  $\mathrm{P} \ne \mathrm{NP}$. Note where the quantum threat falls:
  Shor-type algorithms attack the classical inhabitants of Cryptomania,
  while against the generic hashing of Minicrypt the provable quantum
  gain is only the square root of \cref{sec:barriers}.}
  \label{fig:worlds}
\end{figure}

Two theorems furnish that observation with content, and a third---the
one this section builds toward---turns it into an exact equivalence.

First: \emph{Minicrypt has a canonical witness}. Levin constructed an
explicit, fully specified ``universal one-way function''---a concrete
program one can write down today---which is one-way \emph{if any one-way
function exists at all}~\citep[the tale, and the construction, are told
in][]{Levin2003}. Existence of the whole kingdom is thus equivalent to
the hardness of one specific function, in print since the 1980s. Why,
then, does anyone bother with \SHA{}? Because Levin's function is a
diagonalization---roughly, it runs all short programs and combines
their outputs---and its guarantee is bought at constants and exponents
that make it a philosophical instrument, not an engineering material.
Once any one-way function exists, moreover, the whole private-key
toolbox follows by explicit constructions: H{\aa}stad, Impagliazzo,
Levin, and Luby showed pseudorandom generators can be built from
\emph{any} one-way function~\citep{HILL1999}, with signatures and
message authentication downstream. Read against \cref{sec:gap}: theory
already possesses, with proofs, the entire architecture the world
needs---resting on an unproven ground floor and running too slowly to
use. \SHA{} is the engineering shadow of Levin's function: the object
you deploy when you want the universal guarantee at gigabytes per
second, and are willing to trade the theorem away for it.

Second, the same result read as a boundary: the Akavia--Goldreich--%
Goldwasser--Moshkovitz theorem above says the ground floor cannot be
poured from worst-case \textsc{np}-hardness by standard reductions. The
two faces together frame the epistemic situation exactly: everything
above the one-way-function line is solved mathematics; the line itself
is untethered to the rest of complexity theory.

\subsection{An exact address for the one-way line}\label{sec:liupass}

The third result is recent (2020) and, to our taste, the most
beautiful in the whole subject, because it identifies the one-way-%
function line of \cref{fig:worlds} with the hardness of a problem this
paper has been circling from the start: \emph{detecting structure in
apparently random strings}. Making that precise takes four definitions,
each elementary; we give them carefully, because the payoff is a genuine
``if and only if'' and it is worth being able to read it exactly.

\paragraph{Kolmogorov complexity.} Fix once and for all a universal
Turing machine $U$ (a fixed interpreter: it runs a program $\Pi$ on an
input and prints the result). The \emph{Kolmogorov complexity} of a
string $x \in \bitstrings^{*}$ is the length of the shortest program
that outputs it,
\[
  K(x) \;=\; \min\bigl\{\, |\Pi| \;:\; U(\Pi) = x \,\bigr\}.
\]
This is the honest measure of ``information content'': a patterned
string like $(01)^{n}$ has a tiny generator (``print \texttt{01}
$n$ times''), so $K \approx \log_2 n$, while a structureless string has
no description shorter than itself, $K(x) \approx |x|$. The trouble, for
our purposes, is that $K$ is \emph{uncomputable}---pinning it down
solves the halting problem---so it cannot be the hard problem
cryptography stands on. Something must give the program a deadline.

\paragraph{Time-bounded Kolmogorov complexity.} Fix a polynomial time
bound $t(\cdot)$. The \emph{$t$-time-bounded} complexity restricts the
minimization to programs that print $x$ \emph{within $t(|x|)$ steps}:
\[
  K^{t}(x) \;=\; \min\bigl\{\, |\Pi| \;:\;
      U(\Pi) = x \text{ in at most } t(|x|) \text{ steps} \,\bigr\}.
\]
Now the quantity \emph{is} computable---enumerate all programs of length
$\le |x|$, run each for $t(|x|)$ steps, keep the shortest that prints
$x$---but that brute force costs about $2^{|x|}$. The entire question is
whether one can do better than brute force, and specifically whether one
can do better \emph{on average}, over a random string $x$.

\paragraph{Why almost every string is incompressible.} A counting fact
fixes the landscape (\cref{fig:kolmogorov}). There are only
$2^{n-d}-1 < 2^{n-d}$ programs shorter than $n-d$ bits, so at most a
$2^{-d}$ fraction of the $2^{n}$ strings of length $n$ can have
$K^{t}(x) < n-d$. Compressible strings---those with real structure a
fast program can exploit---are \emph{exponentially rare}; the mass of
$K^{t}$ piles up just below the diagonal $K^{t}(x)=n$. In particular the
\emph{trivial estimate} ``$K^{t}(x) = |x|$''---equivalently, the
one-line algorithm that ignores $x$ and outputs its length---is correct
to within $d$ bits on all but a $2^{-d}$ fraction of inputs. Estimating
$K^{t}$ better than this trivial baseline is precisely the act of
\emph{detecting the rare structure}, and that is the task whose
difficulty the theorem measures.

\begin{figure}[ht]
  \centering
  \begin{tikzpicture}[font=\small]
    % axes
    \draw[->] (0,0) -- (7.4,0) node[below left=0.5mm and -8mm]
      {$K^{t}(x)$};
    \draw[->] (0,0) -- (0,3.5) node[above, align=center, font=\scriptsize]
      {\# of $n$-bit\\ strings};
    \draw[dashed, gray] (6,0) -- (6,3.2) node[above, font=\scriptsize,
      black] {$n$};
    % exponential pile-up near n: bars doubling toward n
    \foreach \k/\h in {1.2/0.10, 1.8/0.14, 2.4/0.2, 3.0/0.28, 3.6/0.42,
                       4.2/0.62, 4.8/0.95, 5.4/1.5, 5.85/2.7} {
      \fill[projframe!75] (\k-0.11,0) rectangle (\k+0.11,\h);
    }
    \node[font=\scriptsize, align=center, projframe] at (2.5,2.35)
      {structured strings:\\ a $2^{-d}$ sliver\\ with $K^{t}<n-d$};
    \draw[->, projframe, font=\scriptsize] (2.5,1.85) -- (3.1,0.45);
    \node[font=\scriptsize, align=center, anchor=east] at (7.3,2.4)
      {almost all mass:\\ $K^{t}(x)\approx n$\\ (incompressible)};
    \draw[decorate, decoration={brace, amplitude=4pt}]
      (0.15,-0.35) -- (4.2,-0.35)
      node[midway, below=4pt, font=\scriptsize]
      {beating ``output $|x|$'' here $=$ detecting structure};
  \end{tikzpicture}
  \caption{The distribution of time-bounded Kolmogorov complexity over
  random $n$-bit strings. By counting, at most a $2^{-d}$ fraction can
  be compressed below $n-d$; the histogram is crushed against the
  incompressible edge $K^{t}\!=\!n$. The one-line ``output $|x|$''
  algorithm is therefore right for nearly everyone---and Liu--Pass say
  that noticeably out-performing it, on average, is exactly as hard as
  inverting every one-way function.}
  \label{fig:kolmogorov}
\end{figure}

\paragraph{Mild average-case hardness, and one-wayness.} Two more
definitions, now the load-bearing ones. Say $K^{t}$ is
\emph{mildly hard-on-average} if some polynomial $p$ makes it
uncomputable on a noticeable slice: for every probabilistic
polynomial-time algorithm $A$ and all large $n$,
\[
  \Pr_{x \,\gets\, \bitstrings^{n}}
     \bigl[\, A(x) \ne K^{t}(x) \,\bigr] \;\ge\; \frac{1}{p(n)} .
\]
The word to dwell on is \emph{mild}: $A$ is permitted to be right on a
$1 - 1/p(n)$ super-majority of strings; hardness on just an
inverse-polynomial sliver suffices. And a \emph{one-way function}---the
one-wayness informally invoked back in \cref{sec:npsearch}---is a
polynomial-time computable $f$ that no probabilistic polynomial-time $A$
can invert except with vanishing probability: for every such $A$, every
polynomial $p$, and all large $n$,
\[
  \Pr_{x \,\gets\, \bitstrings^{n}}
     \bigl[\, f(A(f(x))) = f(x) \,\bigr] \;<\; \frac{1}{p(n)} .
\]
Existence of any such $f$ is the Minicrypt hypothesis---the rung of
\cref{fig:worlds} on which \SHA{}'s conjectured life depends.

\begin{theorem}[Liu--Pass 2020]\label{thm:liupass}
For every polynomial $t(n) \ge (1+\varepsilon)n$ with $\varepsilon>0$
constant, the following are equivalent:
\begin{enumerate}
  \item One-way functions exist (equivalently: secure private-key
  encryption, pseudorandom generators, pseudorandom functions, digital
  signatures, and commitments all exist).
  \item $K^{t}$ is mildly hard-on-average.
\end{enumerate}
Moreover the equivalence is robust to approximation: if one-way
functions exist then no efficient algorithm can even
$(d\log n)$-approximate $K^{t}$ on meaningfully more than the trivial
fraction, for any constant $d$~\citep{LiuPass2020}.
\end{theorem}

The approximation clause is the sharpest way to feel the result. The
do-nothing algorithm ``output $|x|$'' already $(d\log n)$-approximates
$K^{t}$ with probability $\ge 1 - n^{-d}$ (the counting fact again); the
theorem says that beating even \emph{that} baseline, by any
inverse-polynomial margin, is as hard as breaking \emph{every} one-way
function at once. A hair of non-trivial structure-detection, made
reliable, would collapse all of private-key cryptography.

Now read \cref{thm:liupass} in this paper's language. ``$K^{t}(x)$ is
small'' means \emph{$x$ has detectable structure}---a short program
generates it fast. Estimating $K^{t}$ on random inputs is therefore
exactly \emph{hunting for hidden structure in apparently random
strings, on average}---the naturalist's question of \cref{sec:epistemology},
promoted from one function to all strings. The theorem says this hunt is
(mildly) hard \emph{if and only if} objects like \SHA{} can exist at
all. The connection is not a loose analogy; it is the same activity. A
pseudorandom string---say a stretch of the counter stream $\SHA(1),
\SHA(2), \dots$ of \cref{sec:below}---has genuinely \emph{low} $K^{t}$
(the short program ``\,print $\SHA(1..m)$\,'' generates it quickly), yet
is engineered to look maximally incompressible; a distinguisher that
noticed the gap would be estimating $K^{t}$ better than trivial on
exactly these strings, which is to say breaking the generator. So the
empirical program of \cref{sec:gap}---point an instrument at a
pseudorandom object and measure whether any structure pokes through---is
not merely \emph{analogous} to the problem on which the existence of
cryptography rests. Up to the formalization, it \emph{is} that problem,
shrunk to a single function and a laptop. We know no cleaner statement
of why the small experiments of \cref{sec:projects} sit on the main line
of the subject rather than beside it.

\subsection{Two-sided barriers: what is provably out of reach}
\label{sec:barriers}

The deepest results in this corner are \emph{impossibility} theorems,
and they cut in two directions at once.

On the side of \emph{attacks}, the quantum model yields rare
unconditional lower bounds. Grover search inverts any black-box function
on $N$ points in $\Theta(\sqrt{N})$ queries, and this is provably
optimal---no quantum algorithm does better \emph{relative to the
oracle}~\citep{BBBV1997}; for collisions, the Brassard--H{\o}yer--Tapp
birthday-style algorithm costs $\Theta(N^{1/3})$, and Aaronson and Shi
proved a matching $\Omega(N^{1/3})$ quantum lower
bound~\citep{AaronsonShi2004}. These theorems bound what \emph{any}
algorithm can do against a function that behaves as a random oracle, so
they translate into concrete, believed-tight security levels for
\SHA{}---$128$ quantum bits of preimage resistance, and the reason a
$256$-bit digest is considered ``quantum-safe'' for collision resistance.
The catch, familiar by now, is the model: the bounds hold in the
black-box world and become claims about \SHA{} only under the heuristic
that no shortcut exploits its actual circuit.

It is worth pausing to give exponents like $2^{128}$ an empirical
scale, because the planet happens to be running the relevant
experiment. The proof-of-work industry of \cref{sec:uses} evaluates
\SHA{} about $10^{21}$ times per second---roughly $2^{70}$---and its
cumulative total to date is on the order of $2^{96}$ evaluations:
comfortably the largest computation of any kind in human history,
every cycle of it hammering on this one function. Two readings. As a
\emph{statistical} instrument: block arrivals track the uniform model
with textbook fidelity, so this is the most expensive test of the
random-function conjecture ever conducted, passed to date. As a
\emph{yardstick}: the deepest partial inversion all that work has ever
exhibited is a digest beginning with roughly a hundred zero
bits---each mined block is a certified partial preimage, currently
costing about $2^{79}$ evaluations apiece---which still leaves the full
$256$-bit inversion about $2^{156}$ times beyond the species' total
output, and even the $2^{128}$ birthday bound a factor of four billion
past everything computed so far. When \cref{sec:telescope} says the
global geometry of the functional graph is ``operationally
unreachable,'' this is the industrial base against which
``unreachable'' is calibrated.

On the side of \emph{proofs}, the barrier theorems explain the gap of
\cref{sec:gap} as something deeper than a lack of effort. The
\emph{natural proofs} barrier of Razborov and
Rudich~\citep{RazborovRudich1997} shows that any proof of a strong
circuit lower bound which is itself ``constructive and large''---the form
essentially all known lower-bound techniques take---would, ironically,
yield an algorithm breaking the very pseudorandom functions such a lower
bound is meant to certify. In slogan form: \emph{the existence of objects
like \SHA{} is precisely what prevents us from proving that objects like
\SHA{} are hard} by current methods. This is why nobody expects an
unconditional proof of \SHA{}'s security in the foreseeable future, and
why the honest theorems in the subject are either conditional
(\cref{sec:frameworks}), relativized (the quantum bounds above), or about
components rather than the whole (\cref{thm:rivest}).

\subsection{What the proof assistants did prove}\label{sec:verified}

After a section of impossibilities, it is worth recording where
machine-checked mathematics has \emph{succeeded} with \SHA{}---because
it has, completely, one floor down. In 2015 Appel produced a formal,
machine-checked proof, in the Coq proof assistant, that OpenSSL's C
implementation of \SHA{} computes exactly the function specified in
FIPS 180-4: every pointer offset, every rotation, every corner of C
semantics, verified down to a mechanized model of the
compiler's view of the program~\citep{Appel2015}. The HACL\textsuperscript{*}
project carried the method to an entire cryptographic library, written
in the F\textsuperscript{*} proof language and compiled to C; its
verified \SHA{} ships inside Firefox, executing in ordinary users'
browsers~\citep{Zinzindohoue2017}. These are not toy exercises; they
retire a whole species of real historical failure, the
\emph{implementation} bug---and they are proofs about \SHA{} in the
strictest sense of the word.

The epistemological inversion deserves to be savored. \emph{Everything
provable about \SHA{} has now been proved by machines: namely, that the
program is the function.} What no one can prove is any property of the
function itself. Formal verification thus sharpens the boundary of
\cref{sec:gap} into a visible line: below it, mechanized
certainty---code equals specification, checked symbol by symbol; above
it, Rogaway's human ignorance (\cref{sec:definitions}), trust
underwritten by two decades of failed attack. The axiom has quietly
moved. It is no longer ``this function is one-way''; it is ``this
function is the one we wrote down''---and only the second statement is
checkable. A mathematician could not ask for a cleaner exhibit of where,
exactly, in one industrial object, the proved ends and the conjectured
begins.

\subsection{Where this leaves the program}
\label{sec:leaves}

The computability view sharpens, rather than dampens, the case for the
empirical program of this paper. Unconditional hardness is provably out
of reach by known methods; worst-case \textsc{np}-hardness is the wrong
target, and provably unavailable by standard reductions
(\cref{sec:worlds}); the clean theorems that do exist are relativized
bounds assuming random-oracle behavior, or certify the code rather than
the function (\cref{sec:verified}). What is \emph{not} ruled out---by
any barrier theorem---is the discovery of exact, non-statistical
structure in specific reduced systems, of the kind \cref{sec:hope}
catalogued next door and \cref{thm:rivest} exhibited inside the function
itself. The barriers forbid a certain style of \emph{general} proof;
they say nothing against a Josephus-style collapse of a
\emph{particular} finite object. That is exactly the crack the
undergraduate program of \cref{sec:projects} proposes to work in: not to
prove \SHA{} hard---provably hopeless---but to map, instance by instance
and coordinate by coordinate, where the structure ends and the noise
begins.
